\documentclass[11pt]{article}
\usepackage[margin=1.12in]{geometry}
\usepackage{amsmath,amssymb,amsthm,booktabs}
\usepackage{xcolor}
\usepackage[colorlinks=true,allcolors=blue!55!black]{hyperref}
\newtheorem{theorem}{Theorem}
\newtheorem{proposition}[theorem]{Proposition}
\newtheorem{example}[theorem]{Example}
\theoremstyle{definition}
\newtheorem{definition}[theorem]{Definition}
\newcommand{\D}{\mathsf D}
\newcommand{\May}{\mathsf{May}}
\newcommand{\Addr}{\operatorname{addr}}
\newcommand{\push}{_{*}}
\newcommand{\C}{\mathcal C}
\newcommand{\R}{\mathcal R}
\newcommand{\obs}{\mathsf{obs}}
\title{Zero-weight branches in finite probabilistic protocols:\\
  an exact span criterion for sound contextual replacement}
\author{}
\date{Working draft, 20 November 2026}
\begin{document}
\maketitle
\begin{abstract}
Two finite probabilistic protocols can have the same joint law of completed
trees, typed results, and costs, yet disagree on a budget required of every
declared execution. The discrepancy is a zero-weight branch, not a loss of
correlation or a change to the law's observable coordinates. We give a
finite replacement criterion that combines a relation-supported coupling
with coverage of the candidate's declared possibilities. The two conditions
are equivalent to existence of a common weighted history presentation whose
maps preserve probability and tree structure, with a surjective candidate
projection and a directed cost comparison. The criterion composes in a
finite dependent protocol calculus with disjoint ownership and result-dependent
continuation arities. A rational max-flow construction gives local witnesses
or certificates of failure once legal matching pairs are supplied.
The result concerns declared finite executions, not almost-sure safety or
arbitrary schedulers.
\end{abstract}

\paragraph{Provenance.}
AI agents generated most of the underlying mathematical development,
including candidate statements, written proofs and examples, and prepared
this draft. The human project lead directed the research programme and
selected this result for manuscript development. The main arguments are
written proofs, not Lean formalizations. The executable companion checks
finite rational instances of the local criterion under supplied compatibility
data. This draft makes no claim of completed human proof verification,
peer review, or exhaustive priority review.

\section{Equal joint law, unequal all-path budgets}

An annotated finite execution table has \emph{declared} rows with
nonnegative weights summing to one. A zero-weight row is not a stochastic
event to which an execution engine must assign positive frequency. It
describes a branch retained by the specification and counted by its
universal obligations. If an interface treats a row of weight zero as
absent, the distinction studied here disappears at that interface.

Consider a screen with three distinct owned resources: the screen token
itself and two reserve tokens. The screen consumes its token at cost one.
Its low branch returns one continuation hole owning both remaining tokens.
Confirmation uses one token at cost two, and disposal consumes the other
at zero cost. Its high branch returns two different holes, one for each
remaining token. The two assays cost three and five, respectively. They
both run in this declared high branch, without duplicating an owner.
All completed leaves are owner-free. Take the assays deterministic for
this example; adding bit outcomes changes neither the argument nor the
same-law comparison.

\begin{center}
\begin{tabular}{@{}llrrr@{}}
\toprule
Completed tree & Continuation holes after screen &
Cost & Pruned table & Full table\\
\midrule
$\tau_L$ & one bundle & 3 & 1 & 1\\
$\tau_H$ & two disjoint bundles & 9 & absent & 0\\
\bottomrule
\end{tabular}
\end{center}
Both tables have the law $\delta_{(3,\tau_L)}$ on the \emph{full} joint
outcome, including costs and typed results. The pruned table satisfies
an all-path budget of three; the full table does not. Any decision rule
factoring through this law gives the same answer to both tables, irrespective
of its internal representation. Replacing the two completed costs by
integers $1\leq \ell<H$ works with the same owned transitions: charge
$\ell-1$ for confirmation and $H-1,0$ for the two assays after the unit
screen cost.

This also separates two notions of an implementation map. Sending
the one pruned row to the matching full row is a cost-exact native map
in the original model (defined below): the extra target row has an
empty preimage, zero mass and no incoming cost to compare. Its image at
the law level is an identity. But the map is not onto the declared
rows, so this direction of comparison cannot certify that \emph{all}
declarations of the full table are covered.

\section{The finite protocol interface}

\paragraph{Rows and outcomes.}
Write $C=\mathbb N^d$ with coordinatewise addition and order. A fixed
ranked public event signature determines finite planar trees $TX$ with
leaves in $X$ (nullary constructors are allowed). Leaf positions are
addresses, not equivalence classes of equal leaf labels. A full
configuration consists of its public control and the typed private state
of its actual owned tokens. Each continuation leaf is assigned its
declared owned bundle, disjoint from those at the other leaves of the
same tree. Ordinary token names can be carried as data; this does not
confer another owner. All local tables have finite nonempty raw carriers,
normalized weights and finite costs. At each full input a raw table $P$
has a finite set $I$ of distinct row indices, numbers $p_i\geq0$ with
$\sum_i p_i=1$, and outcomes $a_i=(c_i,t_i)\in C\times TX$. A fixed
comparison may erase designated private unary nodes only when its
normalization respects grafting at the corresponding retained addresses;
after any such normalization, the native map equations are asserted on
the normalized presentations, not on the unnormalized body trees.
Retain the private semantic
state needed by subsequent calls. This is not a private-state oracle
for the client.

The law and the possible outcomes are
\[
 \mu_P=\sum_{i\in I}p_i\delta_{a_i},\qquad S_P=\{a_i:i\in I\},
 \qquad \obs(P)=(\mu_P,S_P).
\]
Here $S_P$ is a set of \emph{joint} outcomes, not separate sets of
possible costs and results. In particular $\operatorname{supp}\mu_P
\subseteq S_P$, not necessarily with equality. A predicate $A$ on
joint outcomes holds for all declared rows exactly when $S_P\subseteq A$.
Thus two presentations have the same law and agree on all such predicates
exactly when their pairs agree: a point of the symmetric difference
separates distinct sets. This is an exact observation quotient, not a
faithful representation of row multiplicities or maps of raw rows.
Nor is the whole may-set necessary for scalar upper-bound tests alone:
those see only the maximum cost. Even the law and separate may-marginals
are insufficient for general joint assertions. At law $\delta_{(0,v)}$,
the may-sets
\[
 \{(0,v),(1,t),(2,u)\} \quad\text{and}\quad \{(0,v),(2,t),(1,u)\}
\]
have the same projections onto costs and trees but disagree on the
assertion that every declared outcome with tree $t$ costs at most one.

\paragraph{Native row maps.}
At fixed input, a native map from a table $(I,p,c,t)$ to
$(I',p',c',t')$ over a leaf-state map $v$ consists of a function $r:I\to I'$
such that, for every target row $j$ and source row $i$,
\[
 p'_j=\sum_{r(i)=j}p_i,\qquad
 c'_j\leq\min_{r(i)=j}c_i,\qquad
 t'_{r(i)}=Tv(t_i).
 \tag{1}\label{eq:native}
\]
The minimum is componentwise, with empty minimum $\infty$ in every
coordinate. Equivalently the second condition is the pointwise
inequality $c'_{r(i)}\leq c_i$. A map is cost-exact if equality holds
on each source row; it need not be surjective. A family of these
equations indexed by input pairs gives the native endpoint squares
used here; the input maps select the relevant candidate and reference
fibres.

\paragraph{Related outcomes.}
Use candidate first. Let $P$ and $Q$ be at a related full input pair.
Fix a relation on their typed private configurations that identifies
the compared public interface and matches token identities and kinds
as declared. Put $a\R b$ for outcomes $a=(c,t)$, $b=(d,u)$ when the
compared trees have the same public shape and event labels, their
corresponding leaves have related full configurations, and $c\leq d$.
The first two conditions allow a paired tree over the relation carrier;
they also identify the two sets of addresses. All costs in a certificate
below are taken at the comparison interface, before any unrelated
external erasure.

\section{An exact finite certificate}

\begin{definition}
The local replacement judgment at a related input holds if there is a
coupling $G$ of $\mu_P$ and $\mu_Q$ whose positive entries lie in $\R$,
and for every $x\in S_P$ there is a $y\in S_Q$ with $x\R y$.
The second condition is \emph{candidate coverage}. It does not require
positive mass on its witness or any converse coverage.
\end{definition}

\begin{theorem}[Exact marked-span criterion]\label{thm:span}
The candidate refines the reference if and only if there is a finite
common raw table with two maps satisfying \eqref{eq:native}, to $P$
and $Q$. Require that each common row project to a compatible pair,
have the paired tree over the related configurations, and have exactly
its reference row's cost. Mark the candidate projection as surjective.
The reference projection need not be surjective. These statements hold
at a fixed related input, and the constructions assemble fibrewise
over a finite carrier of related inputs.
\end{theorem}
\begin{proof}
Write $j\in J$ for reference rows, $q_j$ for their weights and $b_j$
for their outcomes, and abbreviate $\mu_P,\mu_Q$ to $\mu,\nu$.
Given the coupling and coverage, first place mass
\[
 \gamma_{ij}=G(a_i,b_j)\frac{p_i}{\mu(a_i)}
                  \frac{q_j}{\nu(b_j)}
\]
on each pair with both denominators positive, and use zero otherwise.
For a positive candidate row, summing the reference factors in any
positive outcome fibre gives one; then the corresponding row sum of
$G$ gives $p_i$. Zero candidate rows have sum zero. The same reasoning
gives the reference marginals. Keep the positive pairs. Each is legal
and every positive candidate row is covered. For each uncovered
candidate row $i$, choose a reference row above a compatible possible
outcome and adjoin that pair with weight zero. Separate candidate rows
require separate lifts even when their outcomes are equal. The
reference partner may itself have zero weight. Neither marginal changes.

Give each common row the paired tree and the reference cost. Equation
\eqref{eq:native} on weights is exactly the two marginal identities;
it holds for costs because compatible candidate cost is no larger
than reference cost, and holds for bodies by projection of paired
leaves. An unmatched reference row is allowed, but its marginal must
be zero. Candidate surjectivity follows by construction.

Conversely push the common weights along the two row maps and then
along the outcome maps. The native marginal equations give a
relation-supported coupling. Lifting a raw representative of any
candidate outcome by the marked surjection gives a declared reference
partner, even when the lifted row has no mass.
\end{proof}

\begin{proposition}[Transfer and exact equality]
Assume candidate coverage. If $x\R y$ and $y\in A_Q$ imply $x\in A_P$,
reference all-path
validity of $A_Q$ implies candidate all-path validity of $A_P$.
In particular this holds for common coordinatewise upper budgets.
At an identical full endpoint, equality of the observation pairs is
equivalent to the existence of a common table with equal joint outcomes
on each pair, cost-exact projections and both projections surjective.
\end{proposition}
\begin{proof}
For transfer, choose a declared reference partner for each declared
candidate row. For equality, if the two observations are equal,
retain every pair of rows with an equal joint outcome $z$. Both
projections are onto. Use pair weight $p_iq_j/\mu(z)$ on positive
outcomes and zero on zero-weight outcomes. This has the required
marginals by summation over each outcome fibre. Conversely the two
surjections give equality of image sets and the two marginal equations
give equality of laws.
\end{proof}

The two conditions in the theorem are independent. The full screen
compared to the pruned screen gives equal laws and no coverage for the
high row. For the other direction use candidates $a,b$ of masses
$3/4,1/4$ and reference rows $a',b'$ of masses $1/4,3/4$. Allow only
matching labels with equal costs, so each candidate row has a partner
but the mass on $a$ cannot be sent to its permitted targets.

\section{Composition at actual continuation addresses}
For a prefix row $i$ with tree $t_i$ and a continuation table at each
actual leaf $s\in\Addr(t_i)$, the composite rows and weights are
\[
 \coprod_i\ \prod_{s\in\Addr(t_i)} I_s,\qquad
 p_i\prod_{s\in\Addr(t_i)}p_{s,k_s}.
 \tag{2}\label{eq:bind}
\]
The composite cost is $c_i+\sum_s c_{s,k_s}$ and the tree is obtained
by grafting each child at its address. The empty tuple has weight one
and adds no cost. Continuations depend only on their full typed leaf
input and admitted framed data, not discarded raw row labels. Leaf
bundles have to be disjoint. If a shared random input is needed across
children it must already be represented in the retained input or
prefix; \eqref{eq:bind} is not a way to construct a missing joint law.
Repeated occurrences with equal labels still make distinct conditional
choices.

For an outcome $z=(c,t)$, let $g_z$ graft and sum a tuple of child
outcomes, and write $(\nu_s,R_s)$ for the child law and may-set at
address $s$. Then the observations of the composite are
\[
 \left(\sum_z \mu(z)(g_z)_*\!\bigotimes_{s\in\Addr(t)}\nu_s,       \bigcup_{z\in S}g_z\!\left[\prod_{s\in\Addr(t)}R_s\right]\right).
 \tag{3}\label{eq:obsbind}
\]
The law uses conditional product followed by mixing, the may-set
uses a Cartesian product followed by union, and both use the same
grafting map. A zero-weight prefix contributes to the union, not the
mixture. To obtain \eqref{eq:obsbind} from \eqref{eq:bind}, first
take the image at each addressed child, then group all prefix rows
with equal full outcome. The child factors are identical within
those groups by the continuation restriction. Relabelling commutes
with grafting; the split between conditional products and mixtures
does not justify multiplying the unconditional marginals.

These formulas are unital and associative. A unit is the single
weight-one zero-cost leaf; this is different from a nullary event,
which receives no children. Both three-stage associations index the
same outer outcome, tuple of middle outcomes and tuple of inner
outcomes at the actual middle addresses. The address bijection for
grafted trees identifies the final tree, the weights are the same
product, and total costs are the same sum. Taking law pushforwards or
images proves the respective equations. Nothing here requires a
nondeterministic set of alternative probability laws.

\begin{theorem}[Closure under finite typed contexts]\label{thm:closure}
Consider finite client terms generated by return, matched public case,
addressed continuation, separated batch on disjoint owned bundles, and
ordinary-data framing and substitution. Acyclic or decreasing-rank
invocation, including finite callbacks, is admitted when its expansion
uses these clauses with the same typed boundaries. Assume:
\begin{enumerate}
\item Initial joint laws have a relation-supported coupling and every declared
candidate initial possibility has a related reference initial possibility.
\item At every common public valuation, every primitive pair (including fixed client primitives) satisfies Theorem~\ref{thm:span}'s criterion
at \emph{every} related full method input, not just those with positive
probability. Case observations, typed interpretations and any public
normalizations are compatible with the compared signatures and grafting.
\item The state relation restricts and reassembles along the designated
disjoint bundles, and related framed ordinary data preserve that relation.
No clause copies a token's ownership or consults a discarded row label.
\end{enumerate}
Then each admitted comparison has the two local conditions on its
composite outcomes and an actual composed native certificate with
surjective candidate projection. For each closing context, all compatible
reference all-path assertions transfer to the candidate.
\end{theorem}
\begin{proof}
Use a local marked span for the prefix and one at each related
addressed child input. Retain a common prefix row with its paired
private states, then choose a tuple of common continuation rows. Give
the tuple the weight, cost and grafting of \eqref{eq:bind}. On fixing
a candidate composite row $(i,(k_s)_s)$, summing the common weights
over lifts of each $k_s$ gives its original child weight at the
candidate's full input. Those marginal sums are independent of which
common prefix row over $i$ was chosen. Summing over the prefix lifts
gives $p_i\prod_s p_{s,k_s}$. The same holds for the reference.
One can also lift any candidate composite row by choosing a prefix
lift, then a child lift at each address irrespective of mass. Empty
tuples are covered. Costs are reference sums and dominate candidate
sums. Related trees graft over the reassembled state relation, so the
native body equations hold and Theorem~\ref{thm:span} applies.

Induct on the well-founded expansion rank (finite acyclic expansions
admit one) and the finite typing derivation at each rank, simultaneously
at all common valuations and related full inputs. Return uses one paired
row, carrying related private entries unchanged. A matched public case uses the
inductive certificate in the selected arm on both sides. A separated
batch uses the same finite-tuple proof conditional on the complete
paired pre-state, with the unused reserve retained unchanged. Ordinary
framing and substitution use the relation preservation assumption.
Public code dispatch has the same tag and equal ordinary, data-only
captures on both sides. Invoke the induction on the callee at strictly
lower rank whenever a call is expanded. Finite acyclic plugging and
callbacks obey the same boundary and rank requirements, including
resource reindexing along typed substitution. Finally apply the raw lift from the proof of
Theorem~\ref{thm:span} to the initial relation (with no tree or cost
comparison needed) and use it as the outer prefix. This preserves any
given initial cross-token correlation. The transfer assertion follows
from coverage.
\end{proof}

For vertical composition at the observation level, use witnesses
$G(x,y),H(y,z)$ with the same middle observation $(\nu,S_Q)$ and compatible full
interfaces. Paste by summing over
$y$ the weights $G(x,y)H(y,z)/\nu(y)$ for $\nu(y)>0$ and zero
otherwise. The row and column sums give the outer laws, and
composing the relations gives compatibility; middle private state
witnesses can be retained before taking images. Compose coverage by
choosing two successive declared partners, even if the middle one
has zero weight. This proves transitivity of the existence judgment
with the composed relation, not canonical interchange of chosen
certificates.

\begin{example}[A nonidentity replacement with zero rows]
At every related method input let the reference split each candidate row
$i$ into an independently supplied finite nonempty audit set $A_i$.
Reference weights are $p_i\alpha_i(a)$ with $\sum_{a\in A_i}\alpha_i(a)=1$;
costs are $c_i+\Delta_{i,a}$, $\Delta_{i,a}\geq0$. Keep the same public
constructor, partition and compared leaf configuration, allowing redundant
related private audit state or designated unary audit nodes the common
normalization removes. Projection to $i$ is onto even if $p_i=0$.
Use its raw indices, weights and costs for the common table at each
related input, with the paired tree over the relation carrier. The
reference map is the identity on indices, and the candidate map is the
indicated projection; both project the normalized paired tree. The map
can decrease cost and identify several positive or zero-weight rows.
By Theorem~\ref{thm:closure}, this local certificate gives all-path
transfer through finite clients satisfying the stated ownership,
initialization and continuation assumptions. No independence of audit
weights from the actual method input or result is needed.
\end{example}

\section{An exact rational checker}
At a single related input, suppose the two raw tables have rational
weights and the legal row pairs have already been determined. Build a
finite network with capacities $p_i$ from the source to each candidate
row, $q_j$ from each reference row to the sink, and capacity one on
allowed candidate--reference edges. A flow of value one is precisely a
compatible raw coupling. When legality is exactly the pullback of $\R$ along the outcome maps,
an outcome coupling lifts to such a flow by Theorem~\ref{thm:span}. For
arbitrary supplied row-level legality, the checker instead solves the
raw-row problem; it does not infer legality from aggregated outcomes. Relative to this requirement, every
positive candidate row is covered. Check each zero candidate row
separately and attach a zero-weight pair to any legal partner, including
a zero-weight reference row when necessary.

If the flow cannot reach one, a minimum cut of capacity less than one
crosses no allowed pair edge, since each such edge alone has capacity
one. For its source-side candidate set $D$, the adjacent reference set
$N(D)$ lies on the source side and $p(D)>q(N(D))$. The inequality
can be checked without rerunning the flow solver. If the zero-row check
fails, the certificate is a declared candidate row with no legal
partner. Conversely either certificate rules out a marked span. For
existence, any cycle in the positive bipartite flow support can be
cancelled by subtracting the smallest alternating edge weight on one
parity and adding it to the other. This leaves both marginals unchanged
and deletes at least one positive edge. Repeating gives a forest on
the two positive supports and a witness with at most
$|I_+|+|J_+|-1$ positive rows. Both supports are nonempty by normalization.
The zero fillers give a total bound of $|I|+|J_+|-1$ rows.

The companion script works with exact rational arithmetic, accepts
the relation as a supplied finite set, and reports either kind of
failure or a witness checked against the two marginals and coverage.
It does not type-check the language, find the intended relation, decide
closure for all contexts, or serve as a Lean proof.

\section{Scope and relation to prior work}
Support relations have a long history in probabilistic simulation:
the finite weight-function formulation of Segala and Lynch requires a
matching relation on positive entries \cite{segalalynch}. Probabilistic
product programs can express coupling witnesses with the original
marginals \cite{barthe}. Our mass condition is of this familiar kind.
It does not enforce coverage of a declared row that carries no mass.
The additional condition here is exact for a specific certificate:
two original directed row maps, with surjectivity on the candidate
side. It is not an assertion that an arbitrary program admits an
executable product program or that contextual validity implies a span.

Varacca and Winskel's indexed valuations retain positive
multiplicity while identifying zero indices; their goal is to combine
probabilistic and nondeterministic choice by a distributive law
\cite[Defs.~4.1--4.5]{varaccawinskel}. Here there is one fixed law, a
separate may-set, and no nondeterministic choice among laws. Row
multiplicity is discarded in the observation and retained in its native
certificate. In particular, the construction does not give a
distributive law between the ordinary finite distribution and powerset
monads for arbitrarily combining these effects.

Outcome Logic defines support relative to the chosen weight semiring
\cite[Def.~3.6 and Sec.~6.1]{outcomelogic}. Ordinary probability weight
zero does not contribute to that support. This is different from
claiming that semiring-parametric reasoning cannot see a declared
zero-weight row: pairs $(p,b)$ where $b$ is a presence bit and $p>0$
implies $b=1$ give a concrete way to encode it. Addition uses
$(+,\lor)$ and multiplication uses $(\cdot,\land)$; $(0,1)$ is then
not the additive zero. Normalized finite distributions in those
weights recover the law and may-set. We use the elementary
bookkeeping directly and do not claim to instantiate the general proof
system of Outcome Logic. The native cost-lowering projections and
one-sided raw coverage are the additional certificate structure.

The theorem does not say that a single direct native map must exist.
Nor does allowing marked spans make the old category faithful under
the observation quotient: duplicate equal outcomes and distinct row
maps can still be identified. The proof uses finite independent
choices \emph{conditional on the retained prefix}; changing to joint
new draws across holes or allowing one continuation to see hidden
raw identities changes the composition law. Nor does the result
extend to arbitrary recursion or almost-sure semantics by replacing
finite sums with limits. The present comparison intentionally
preserves the distinction between declared and measure-supported paths.

\paragraph{Source and evidence.}
The manuscript develops a finite publication slice of the project
dossier \href{https://github.com/eternal-november/eos-research-programmes/blob/f458186590df968a516ba73a8eb76fcf044544b0/projects/contextual-composition/protocols/law-and-possibility/Law_and_Possibility_Repair_for_Dependent_Protocols.tex}{\emph{Law and Possibility Repair for Dependent Protocols}}
(especially the sections on the exact marked original span and
contextual closure). The original calculus supplies the native map
equations and finite typed protocol operations; the earlier
refinement development supplies conditional coupling induction and the
raw-lift formula. The addition emphasized here is exact recognition
by a candidate-surjective history span and its all-path use in the
same occurrence calculus, not the novelty of conditional couplings or
of distinguishing distributions from sets. The project dossier also
contains an enriched coefficient presentation and a restricted
reconstruction comparison; neither is needed for the statements above.
\begin{thebibliography}{9}
  \raggedright
\bibitem{segalalynch}
Roberto Segala and Nancy A. Lynch.
Probabilistic simulations for probabilistic processes.
In \emph{CONCUR '94}, LNCS 836, pp.~481--496. Springer, 1994.
\href{https://doi.org/10.1007/978-3-540-48654-1_35}{doi:10.1007/978-3-540-48654-1\_35}.

\bibitem{barthe}
Gilles Barthe, Benjamin Gr\'egoire, Justin Hsu, and Pierre-Yves Strub.
Coupling proofs are probabilistic product programs.
In \emph{POPL '17}, pp.~161--174. ACM, 2017.
\href{https://doi.org/10.1145/3009837.3009896}{doi:10.1145/3009837.3009896}.

\bibitem{varaccawinskel}
Daniele Varacca and Glynn Winskel.
Distributing probability over non-determinism.
\emph{Mathematical Structures in Computer Science}, 16(1):87--113, 2006.
\href{https://doi.org/10.1017/S0960129505005074}{doi:10.1017/S0960129505005074}.

\bibitem{outcomelogic}
Noam Zilberstein.
Outcome Logic: A unified approach to the metatheory of program logics with
branching effects.
\emph{ACM Transactions on Programming Languages and Systems}, 47(3),
article 14, 2025.
\href{https://doi.org/10.1145/3743131}{doi:10.1145/3743131}.
\end{thebibliography}
\end{document}
